\documentclass[11pt,letterpaper]{article}

\usepackage[T1]{fontenc}
\usepackage{amsmath,amssymb}
\usepackage{tikz}
\usepackage{pgfplots}
\pgfplotsset{compat=1.18}
\usepackage{algorithm}
\usepackage{algpseudocode}
\usepackage{booktabs}
\usepackage{enumitem}
\usepackage[hidelinks]{hyperref}

\newcommand{\N}{\mathbb{N}}
\newcommand{\Z}{\mathbb{Z}}
\newcommand{\Q}{\mathbb{Q}}
\newcommand{\R}{\mathbb{R}}
\newcommand{\C}{\mathbb{C}}
\newcommand{\Map}[3]{#1 \colon #2 \to #3}
\newcommand{\Mod}[1]{#1\text{-}\mathrm{Mod}}
\newcommand{\GroupHom}{\mathrm{Hom}_{\mathbf{Grp}}}
\newcommand{\RingHom}{\mathrm{Hom}_{\mathbf{Ring}}}

\newtheorem{theorem}{Theorem}
\newtheorem{lemma}[theorem]{Lemma}
\newtheorem{corollary}[theorem]{Corollary}
\newtheorem{definition}[theorem]{Definition}
\newtheorem{remark}[theorem]{Remark}
\newenvironment{proof}{\noindent\textit{Proof.}}{\hfill\(\square\)}

\title{Division-Time Laws at the Full Residue-Cycle Boundary}
\author{}
\date{\today}

\begin{document}
\maketitle

\begin{abstract}
When the affine residue cycle has length equal to the divisor, division never leaks away: every orbit has infinitely many divisions. The first-division return nevertheless has an exact uniform word law. We prove that its branch codes are the full one-sided shift, derive density-one strict descent when the mean logarithmic multiplier is negative, and give explicit classical bounds with constants \(243/256\) and \(65536/84375\). A separate common-divisor obstruction classifies a non-coprime family by an exact affine tail. All statements concern the division-first raw map.
\end{abstract}

\section{First-division returns}

Fix integers \(q\ge2\), \(c\ge1\), and \(D\ge2\), and define
\[
T(n)=
\begin{cases}
n/D,&D\mid n,\\
qn+c,&D\nmid n.
\end{cases}
\]
Assume initially that \(\gcd(q,D)=1\). The affine permutation
\(F(r)=qr+c\pmod D\) has a least positive return \(m\) of \(0\).
This article treats \(m=D\), so its zero orbit contains every residue.
For each state \(n\), let \(r(n)\) be the least nonnegative number of affine
steps needed to reach a multiple of \(D\).
Then \(0\le r(n)\le D-1\), and its first-division return is
\[
H(n)=\frac{q^{r(n)}n+cS_{r(n)}}{D},
\qquad
S_r=\sum_{j=0}^{r-1}q^j,\quad S_0=0.
\]
The raw time of this return is \(r(n)+1\), not \(r(n)\).
Every positive seed stays positive, and every return exists.

\begin{lemma}[Exact marked cylinders]
\label{lem:cylinders}
For every \(h\ge0\) and every word
\(w=(r_0,\ldots,r_{h-1})\in\{0,\ldots,D-1\}^h\), the positive integers
with that first-return word form exactly one residue class modulo \(D^h\).
On that class,
\[
H^h(n)=A_w n+C_w,\qquad
A_w=\frac{q^{R_w}}{D^h},\qquad R_w=\sum_{j<h}r_j,
\]
where \(C_w\ge0\). Every word occurs, and distinct words give distinct classes.
\end{lemma}

\begin{proof}
A residue with first-hit time \(r\) is the unique inverse image \(F^{-r}(0)\).
It has no earlier hit: otherwise the zero orbit would have a positive
return smaller than \(D\). These \(D\) residues are distinct.

Suppose a length-\(h\) word has representative \(a\), \(0\le a<D^h\),
endpoint \(E=H^h(a)\), and numerator multiplier \(P=q^{R_w}\).
For every \(t\ge0\), states \(a\) and \(a+D^h t\) have the same preceding
word and
\[
H^h(a+D^h t)=E+Pt.
\]
To justify this assertion at each preceding block, the incoming difference
is divisible by the remaining power of \(D\). Affine multiplication by \(q\)
preserves it and division removes exactly one power. Thus the current
residues, and therefore their first-hit times, agree. Applying the affine
return formula gives the displayed endpoint difference.

Since \(\gcd(P,D)=1\), the residues \(E+Pt\), \(0\le t<D\), run through all
residues modulo \(D\). Each desired next first-hit time therefore selects
exactly one high digit \(t\). This proves existence and uniqueness of
each child class. Child classes from different parent classes cannot
coincide modulo \(D^h\). Starting with the sole depth-zero class proves
the result by induction. Composing the return formulas gives \(A_w\)
and a nonnegative additive term \(C_w\).
\end{proof}

\section{The exact random law and its scope}

\begin{corollary}
\label{cor:moments}
Uniform counting of residues modulo \(D^h\) makes the \(h\) return times
independent uniform variables on \(\{0,\ldots,D-1\}\). In particular,
\[
\mu=\frac{D-1}{2},\qquad
\sigma^2=\frac{D^2-1}{12},\qquad
\mathbb E R_w=h\mu,\qquad
\operatorname{Var}(R_w)=h\sigma^2.
\]
\end{corollary}

\begin{proof}
Lemma~\ref{lem:cylinders} is a bijection between the \(D^h\) residues and
the \(D^h\) words. A prescribed set of \(k\) positions therefore has
\(D^{h-k}\) completions. This proves the product law. The formulas for
one uniform digit follow from
\(\sum_{r<D}r=D(D-1)/2\) and
\(\sum_{r<D}r^2=D(D-1)(2D-1)/6\).
Cross terms in the variance vanish under the proved product law.
\end{proof}

Let \(\Z_D=\varprojlim_h\Z/D^h\Z\), with its uniform probability measure.
This definition allows composite \(D\). Division by \(D\) is defined only
on its image, so it does not replace the raw guard by field divisibility.

\begin{theorem}[Shift coding]
\label{thm:shift}
The first-return word gives a homeomorphism
\(\Z_D\cong\{0,\ldots,D-1\}^{\N}\) conjugating \(H\) with deletion of the
first symbol. It carries uniform measure to the uniform product measure.
\end{theorem}

\begin{proof}
The compatible marked cylinders in Lemma~\ref{lem:cylinders} give one
residue modulo \(D^h\) for every prefix of length \(h\). An infinite word
therefore determines a unique element of the inverse limit. Conversely
an inverse-limit element determines every prefix. Cylinders are bases
for both topologies, proving continuity in both directions. Applying
\(H\) removes the first return block, so the codes intertwine the two
maps. A cylinder of length \(h\) has measure \(D^{-h}\) on both sides,
which proves equality of the measures.
\end{proof}

The real-valued linear cocycle is
\[
\lambda(x)=r(x)\log q-\log D,\qquad
\gamma=\int\lambda\,dx=\frac{D-1}{2}\log q-\log D.
\]
It is bounded and defined from the branch marks, not from an independent
resampling of an integer trajectory. An integrable change of coordinates
replaces it by \(\lambda+\psi\circ H-\psi\), whose integral is still
\(\gamma\), because \(H\) preserves uniform measure. Thus \(\gamma\ne0\)
obstructs an integrable coboundary. Periodic marked words retain their
own sums; averaging does not replace their closing equations.

The shift coding establishes a precise connection with a Bernoulli
process. It does not make every positive integer a typical inverse-limit
point. The countable embedded set of integers has uniform measure zero.
The integer conclusions below are obtained instead from finite cylinder
counts and explicit bounds on additive terms.

\begin{lemma}[Contracting full-cycle parameters]
\label{lem:parameters}
Under the standing hypotheses \(q\ge2,c\ge1,D\ge2\),
\(\gcd(q,D)=1\), and \(m=D\), one has \(\gamma<0\) if and only if
\[
q=3,\qquad D=2,\qquad c\ \text{is odd}.
\]
\end{lemma}

\begin{proof}
The affine permutation has one cycle containing every residue, so it has no fixed point when \(D\ge2\).
If \(q=2\), the residue \(-c\) is a fixed point:
\(2(-c)+c=-c\).
Thus full access excludes \(q=2\), and \(q\ge3\).

For \(D\ge3\), induction gives \(3^{D-1}\ge D^2\).
At \(D=3\) both sides equal \(9\).
For the induction step,
\[
3D^2-(D+1)^2=2D^2-2D-1>0\qquad(D\ge3),
\]
so multiplying the preceding inequality by three proves the next one.
Hence \(q^{D-1}\ge3^{D-1}\ge D^2\).
Taking logarithms shows \(\gamma\ge0\).
Negative drift therefore forces \(D=2\).
Its inequality is then \(q<4\), and \(q\ge3\) forces \(q=3\).
Modulo two the affine map is \(r\mapsto r+c\), whose zero cycle has
length two exactly when \(c\) is odd.

Conversely these parameters give full access and
\(2\gamma=\log3-\log4<0\).
The arithmetic bounds used here are formalized as
\texttt{ThreeSquareGrowth} and \texttt{NegativeGainParameters};
the exclusion of a full cycle with \(q=2\) is the fixed-point argument above.
\end{proof}

Thus the negative-drift conclusion below applies to all positive odd-offset dyadic maps \(F_{3,c,2}\), not to additional positive full-cycle multiplier--divisor pairs.
The exact full-shift law still holds when the drift is nonnegative.

\section{Closing descent on positive integers}

Set
\[
A_*=\frac{q^{D-1}}D,\qquad
C_*=\frac{cS_{D-1}}D,\qquad
M_h=C_*\sum_{j=0}^{h-1}A_*^j.
\]
Every word has \(C_w\le M_h\): each return has nonnegative additive term
at most \(C_*\) and multiplier at most \(A_*\), so induction on the
composition length proves this bound.

\begin{theorem}[Density-one strict descent]
\label{thm:descent}
If \(\gamma<0\), the set of positive integers \(n\) for which
\(T^k(n)<n\) at some positive raw time \(k\) has ordinary natural density
one.
More explicitly, choose \(\eta>0\) such that
\(\rho=q^{\mu+\eta}/D<1\), and for \(h\ge1\) put
\[
N_h=\left\lfloor\frac{M_h}{1-\rho^h}\right\rfloor+1.
\]
For every integer \(X\ge1\), the number \(E(X)\) of positive integers
\(n\le X\) with no strict descent at any positive raw time satisfies
\[
E(X)\le N_h+\frac{\sigma^2}{h\eta^2}(X+D^h).
\]
\end{theorem}

\begin{proof}
The choice of \(\eta\) is possible because
\(\mu\log q<\log D\) and \(\log q>0\).
Call a word good when \(R_w\le h(\mu+\eta)\). On a good class,
\[
H^h(n)\le \rho^h n+M_h<n
\quad\hbox{whenever }n\ge N_h.
\]
The corresponding raw time is the sum of \(h\) positive block lengths,
and is at most \(hD\).

For the uniform finite word space, every bad word has
\((R_w-h\mu)^2>h^2\eta^2\). Summing this nonnegative square over the bad
words and then over all words proves, without a limiting probability
argument, that their proportion is at most \(\sigma^2/(h\eta^2)\).
There are therefore at most
\(D^h\sigma^2/(h\eta^2)\) bad classes. Each class has at most
\(X/D^h+1\) representatives among \(1,\ldots,X\).
All non-descending seeds above \(N_h\) are in these classes, proving
the displayed estimate.

For fixed \(h\), divide by \(X\) and let \(X\) tend to infinity.
The upper limit is at most \(\sigma^2/(h\eta^2)\).
Now let \(h\) tend to infinity. The upper limit is zero, proving
ordinary natural density one. This order of limits fixes each finite
additive bound before it is used.
\end{proof}

\begin{remark}
Strict descent at one time is not eventual arrival at a particular cycle.
The theorem does not remove the density-zero exceptional set or assert
that repeatedly selected smaller seeds satisfy the same typicality
condition.
\end{remark}

\section{Explicit classical constants}

For \((q,c,D)=(3,1,2)\), a return time is \(0\) or \(1\), and
\[
H(n)=
\begin{cases}
n/2,&n\text{ even},\\
(3n+1)/2,&n\text{ odd}.
\end{cases}
\]
For \(h=8t\), call a word good if its number \(a\) of ones is at most
\(5t\). Its multiplier is at most
\[
\frac{3^{5t}}{2^{8t}}=\left(\frac{243}{256}\right)^t.
\]
Let \(\beta=65536/84375\), so \(0<\beta<1\).

\begin{lemma}[Tilted binary count]
\label{lem:binary}
The proportion of bad words of length \(8t\) is at most \(\beta^t\).
\end{lemma}

\begin{proof}
Give a word with \(a\) ones weight \(5^a3^{8t-a}\).
The sum of all weights is \((5+3)^{8t}=8^{8t}\).
A bad word has \(a>5t\), and replacing a zero by a one increases its
weight, so every bad weight is at least \(5^{5t}3^{3t}\).
Consequently
\[
\frac{\#\{\hbox{bad words}\}}{2^{8t}}
\le
\frac{8^{8t}}{2^{8t}5^{5t}3^{3t}}
=
\left(\frac{4^8}{5^5\,3^3}\right)^t
=\beta^t.
\]
These are exact finite counts.
\end{proof}

In the classical specialization choose \(\eta=1/8\), so that
\(\rho^{8t}=(243/256)^t\) in the definition of \(N_h\).
Here \(M_h\le(3/2)^h-1\). For \(t\ge1\),
\[
N_{8t}\le
\frac{(3/2)^{8t}}{1-243/256}+1
\le 22\cdot256^t.
\]
The second inequality follows from
\((3/2)^8<256\) and \(256/13+1<22\).
Thus
\[
E(X)\le22\cdot256^t+\beta^t(X+256^t).
\]

\begin{corollary}[A quantitative exceptional-set bound]
\label{cor:power}
Put
\[
\alpha=\frac{\log(1/\beta)}{\log(256/\beta)}\in(0,1).
\]
For every integer \(X\ge256/\beta\),
\[
E(X)\le (22+2/\beta)\,X^{1-\alpha}.
\]
\end{corollary}

\begin{proof}
Choose \(t=\lfloor\log X/\log(256/\beta)\rfloor\), which is positive.
Then \(256^t\le X^{1-\alpha}\) and
\(\beta^t\le\beta^{-1}X^{-\alpha}\).
Also \(256^t\le X\). Substitute these three bounds into the preceding
finite estimate.
\end{proof}

% Two quarter-budget bounds close fixed-fraction descent with the raw logarithmic clock.
\begin{theorem}[Fixed-fraction descent in logarithmic raw time]
\label{thm:fractional}
Fix a positive odd integer \(c\) and a real \(C\ge1\), and use the
raw map \(T=F_{3,c,2}\). Put
\[
\begin{aligned}
\rho&=243/256, & G&=(3/2)^8, & L&=\log(256/\beta),\\
t_C&=\max\left(1,\left\lceil\frac{\log(4C)}{\log(1/\rho)}\right\rceil\right).
\end{aligned}
\]
Let \(\mathcal E_{c,C}(X)\) count the positive integers \(n\le X\)
for which there is no raw time
\[
0\le j\le\frac{16\log n}{\log G}\quad\text{with}\quad T^j(n)<n/C.
\]
For every integer \(X\ge\exp(Lt_C)\),
\[
\mathcal E_{c,C}(X)
\le(4Cc+1+2/\beta)X^{1-\alpha},
\qquad \alpha=\frac{\log(1/\beta)}L>0.
\]
In particular, for every fixed \(c,C\), density-one seeds reach a
value less than their starting value divided by \(C\) in a logarithmic
number of raw ticks. The bound is ordinary natural density, with
explicit fixed-parameter constants.
\end{theorem}

\begin{proof}
The depth-\(h\) offset for \((3,c,2)\) obeys
\(C_{j+1}\le(3/2)C_j+c/2\), \(C_0=0\), because a return has delay
zero or one. Induction gives
\(C_h\le c((3/2)^h-1)\). For a good depth-\(8t\) word, therefore,
\[
H^{8t}(n)\le\rho^t n+cG^t.
\]
If \(t\ge t_C\), then \(\rho^t\le1/(4C)\). If also
\(n\ge4CcG^t\), the offset is at most \(n/(4C)\). Hence
\[
H^{8t}(n)\le n/(2C)<n/C.
\]
Each first-division block takes at most two raw ticks, so this endpoint
is reached at a raw time \(j\le16t\). Since \(C,c\ge1\), the same
lower bound on \(n\) gives \(n\ge G^t\), and consequently
\(j\le16\log n/\log G\). Both the endpoint bound and the clock bound
hold for the same legal word and seed.

Exactly one class modulo \(256^t\) realizes each marked word.
Lemma~\ref{lem:binary} bounds the fraction of bad classes by \(\beta^t\).
Each class contributes at most \(X/256^t+1\) seeds to \([1,X]\).
Discarding the at most \(\lceil4CcG^t\rceil\) small seeds gives
\[
\mathcal E_{c,C}(X)
\le 4CcG^t+1+\beta^t(X+256^t).
\]
Choose \(t=\lfloor\log X/L\rfloor\). The assumed threshold ensures
\(t\ge t_C\). As in Corollary~\ref{cor:power},
\[
256^t\le X^{1-\alpha}\le X,
\qquad \beta^t\le\beta^{-1}X^{-\alpha}.
\]
Since \(1<G<256\), also \(G^t\le X^{1-\alpha}\), and
\(1\le X^{1-\alpha}\). Substitution proves the stated estimate.
The exceptional fraction tends to zero since \(\alpha>0\).
\end{proof}

This theorem controls one fixed-factor stopping event, not capture of a
fixed terminal value. The constants depend on \(c,C\), and the same
counting argument cannot be reapplied along an adaptively selected
trajectory without additional control of its selected descendants.

At depth twelve the exact contracting-word counts are \(3302/4096\)
for \(q=3\) and \(1586/4096\) for \(q=5\), both with \(c=1,D=2\).
Both models have all \(4096\) marked words, with delay sum \(24576\)
and squared-delay sum \(159744\). Their differing contraction counts
come from the multiplier, not from a different word law.
The computations in \cite{BoundaryCode26} enumerate words and check
these integer identities; they do not test an infinite predicate.

\section{A non-coprime classification}

The affine-escape mechanism below permits non-coprime parameters.
It is also proved in the module
\nolinkurl{QCD.AffineEscape.001/lean/AffineEscape.lean}.

\begin{theorem}[Common-divisor obstruction]
\label{thm:noncoprime}
Let \(q\ge1,c\ge1,D\ge2\). Suppose that an integer \(s\ge1\) divides
both \(q\) and \(D\), but does not divide \(c\).
Every positive seed has a unique representation \(n=D^k u\), where
\(k\ge0\), \(u>0\), and \(D\nmid u\). Its raw orbit satisfies
\[
T^{k+j}(n)=q^j u+cS_j\qquad(j\ge0),
\]
and tends to infinity along the full raw sequence.
In particular the conclusion holds when \(\gcd(q,D)\nmid c\).
\end{theorem}

\begin{proof}
If \(D\mid qx+c\), then \(s\mid qx+c\). Since \(s\mid qx\), this
would imply \(s\mid c\), a contradiction. Thus every affine image
is permanently outside the division branch.

Extract factors of \(D\) from \(n\). Each legal quotient is a smaller
positive integer because \(D\ge2\), so extraction terminates at \(u\).
If \(D^k u=D^\ell v\) with both remainders indivisible by \(D\), unequal
exponents would make one remainder divisible by \(D\); equal exponents
give \(u=v\) by cancellation. The first \(k\) raw events are therefore
exact divisions and give \(u\). Thereafter every event is affine,
and induction proves the geometric-sum formula. Since \(q,c\ge1\),
each affine event increases the state by at least one. Removing the
finite prefix proves full-sequence escape.
\end{proof}

For example, \((q,c,D)=(6,1,4)\) has this property with \(s=2\).
The hypothesis matters: \((2,2,2)\) instead has the positive raw cycle
\(1\mapsto4\mapsto2\mapsto1\).
The case \(\gcd(q,D)\mid c\) requires different branch information.

\section{Ordered affine structure}

Each legal raw word is an ordered product of the two affine matrices
\[
M_{\mathrm{aff}}=
\begin{pmatrix}q&c\\0&1\end{pmatrix},
\qquad
M_{\mathrm{div}}=
\begin{pmatrix}1/D&0\\0&1\end{pmatrix}.
\]
Its product acts as \(x\mapsto Ax+B\).
The coefficient \(A=q^a/D^b\) records multiplication and division
counts; \(B\) records the ordered offset insertions.
A rational matrix product does not certify integrality or legality:
these are the original divisibility and nondivisibility conditions
at every intermediate state.

This structure explains several different proved mechanisms.
At \(c=0\), the additive coordinate vanishes and multiplicative
dependence controls scalar lattice trajectories. At \(q=1\),
congruence access and additive contraction control the return.
For \(m<D\), loss of legal cylinder children controls density.
At \(m=D\), no children are lost, but the marked multiplier law
can still close a descent estimate. The common-divisor theorem
uses an arithmetic obstruction before any averaging.
These are distinct conclusions obtained from preserved parts
of one ordered composition, not a complete orbit invariant.

\section{Formal and computational companions}

The modules \texttt{BoundaryDrift.lean} and \texttt{NonCoprime.lean}
prove the tilted binary inequalities, closing constants, the binary
offset and endpoint envelopes, the quarter-budget fixed-fraction
closing estimate, and common-divisor affine-tail classification.
The complete fixed-fraction counting theorem and logarithmic raw-clock
bound above are written proofs, not Lean declarations. The existing
\texttt{ResidueCylinders.lean} proves legal source lifting, exact
raw-clock recovery, and infinitely many divisions at \(m=D\).
The marked shift conjugacy and density transfer in
Theorems~\ref{thm:shift} and \ref{thm:descent} are proved in this article;
they are not asserted as declarations in those Lean modules.

\bibliographystyle{amsalpha}
\bibliography{BoundaryReferences}
\end{document}
